\documentclass[uplatex,dvipdfmx,11pt,a4paper]{jsarticle}

\usepackage[a4paper,margin=24mm]{geometry}
\usepackage[dvipdfmx]{graphicx}
\usepackage{booktabs}
\usepackage{longtable}
\usepackage{amsmath}
\usepackage{caption}
\usepackage{float}
\usepackage{xcolor}
\usepackage{url}
\usepackage[dvipdfmx]{hyperref}
\usepackage{pxjahyper}

\hypersetup{
  colorlinks=true,
  linkcolor=blue,
  urlcolor=blue,
  citecolor=blue,
  pdftitle={法然・親鸞の共有核とはみ出し領域：三層探索地図による『選択集』・『教行信証』比較},
  pdfauthor={上山大信},
  pdfkeywords={法然, 親鸞, 選択本願念仏集, 教行信証, 浄土教, 共有核, はみ出し領域, 意味埋め込み, 三層分析, SAT, text-embedding-3-large}
}

\captionsetup{font=small,labelfont=bf}
\renewcommand{\figurename}{図}
\renewcommand{\tablename}{表}
\renewcommand{\refname}{参考文献}
\renewcommand{\abstractname}{要旨}

\title{法然・親鸞の共有核とはみ出し領域：\\[0.35em]三層探索地図による『選択集』・『教行信証』比較}
\author{上山大信}
\date{2026年6月5日}

\begin{document}

\maketitle

\begin{abstract}
本稿は、法然『選擇本願念佛集』と親鸞『顯淨土眞實教行證文類』が、意味空間上でどのように重なり、また共有核から相対的に離れて現れる「はみ出し領域」（相対的突出領域）がどの方向に現れるかを、意味層・文体語彙層・典拠マーカー層の三層から検討する探索的研究である。文体語彙層は厳密な文体計量ではなく、辞書ベースの語彙・主題濃淡として扱う。筆者は先に「意味・文体・典拠マーカーの三層地図による仏教文献探索：宗派別参照群・阿弥陀経二訳・親鸞文献の予備的分析」において、宗派別参照テキスト群、阿弥陀経二訳、親鸞文献を対象に三層地図を試みた\cite{ueyama2026three-layer-map}。本稿ではこの枠組みを法然・親鸞・浄土教祖師文献アンカーへ限定し、両者が重なるかどうかではなく、重なった後に何が分岐するかを問う。『選択集』と『教行信証』の比較は従来から教学史・思想史の文脈で論じられてきたが、本稿はその関係を新たに結論づけるものではなく、従来研究が扱ってきた継承・展開の問題を、本文全体の近傍構造と三層分析によって探索的に可視化することを目的とする。

結果として、法然チャンク69件の最近傍非自己は親鸞50件、道綽11件、源信4件、善導3件、曇鸞1件であり、法然チャンクの多くは、親鸞文献チャンクと意味空間上で強く重なることが確認された。一方、親鸞チャンク191件の最近傍非自己は法然48件、道綽45件、源信36件、曇鸞23件、天親21件、善導17件、龍樹1件に分かれた。参照群のチャンク数バイアスを確認するため、各参照群の候補数を上限20件に制限した1000回のサブサンプリングも行ったところ、法然から親鸞への近接比率は低下するが、親鸞・道綽を中心とする傾向は残り、親鸞側では道綽・法然・天親・曇鸞などへの分散が確認された。

共有核からのはみ出し領域を比較すると、法然側では正雑二行、諸行との取捨、本願、付属・証誠など、念仏を諸行から選び出す論証が現れた。これに対し親鸞側では、信巻の信・三心・罪救済、化身土巻の護法・鬼神・宇宙秩序、外教批判、方便・真仮整理が目立った。したがって本稿の仮説は、法然と親鸞の差は念仏そのものの意味差ではなく、同じ念仏圏をどのような論証形式・典拠体系へ展開するかの差として現れる、というものである。分析にはSAT漢文本文、Unicode-safeな700トークンチャンク、100トークン重なり、OpenAI \texttt{text-embedding-3-large} による3072次元埋め込みを用いた。ただし、この仮説は埋め込み、近傍関係、語群カウント、典拠マーカーから得た探索的結果であり、個別箇所について既存の教学史・引用研究と全面的に照合した結論ではない。引用文の立場と著者自身の立場の分離、巻別・章別境界の精密化、古写経・異本研究との照合は今後の課題である。
\end{abstract}

\noindent\textbf{キーワード：} 法然，親鸞，選択本願念仏集，教行信証，共有核，はみ出し領域，意味埋め込み，三層分析，SAT，典拠マーカー

\section{はじめに}

\subsection{本稿の位置づけ}

本稿は、拙稿「意味・文体・典拠マーカーの三層地図による仏教文献探索：宗派別参照群・阿弥陀経二訳・親鸞文献の予備的分析」\cite{ueyama2026three-layer-map}を受け、対象を浄土教文献に狭めて再検討する探索的研究である。前稿では、漢訳仏典および日本撰述仏教文献を、意味層、文体・語彙層、典拠マーカー層に分けて可視化する枠組みを試みた。そこでは、意味的に近いこと、文体的に近いこと、典拠・引用経路として近いことは同一ではない、という点を中心命題とした。

本稿では、対象をより狭く取り、浄土教、とくに法然と親鸞の関係に集中する。法然の代表的著作として『選擇本願念佛集』を、親鸞の代表的著作として『顯淨土眞實教行證文類』を中心対象に据える。以下、必要に応じて前者を『選択集』、後者を『教行信証』と略記する。両書は、法然・親鸞の浄土教思想を比較するうえで中心的資料であるため、本稿の中心対象とする。両者は浄土教の中心問題を共有するため、意味空間上でも大きく重なることが予想される。したがって、本稿で重要なのは「重なるかどうか」ではなく、重なった上で、どの方向に差異が出るかである。

本稿では、この問いに対し、SAT大正新脩大藏經テキストデータベースから取得した漢文本文を用い、法然・親鸞・浄土教祖師文献アンカーの探索地図を作成する。さらに、共有核、法然側のはみ出し領域、親鸞側のはみ出し領域を、意味層、文体語彙層、典拠マーカー層の基本構成で読む。ここでいう「はみ出し領域」とは、共有核から相対的に離れて現れる本文領域、すなわち相対的突出領域を指す探索的な作業概念である。

\subsection{先行研究と本稿の位置づけ}

法然『選擇本願念佛集』と親鸞『顯淨土眞實教行證文類』の関係については、従来から教学史・思想史・引用研究の立場から論じられてきた。森脇一掬「選択集と教行信証に関する一考察」\cite{moriwaki1953}、稲葉秀賢「『教行信証』と『選択集』」\cite{inaba1972}、梅原隆章「選択集と教行信証」\cite{umehara1974}、浅井成海「『選択集』と『教行信証』――基本的問題の継承と展開――」\cite{asai2004}などは、両書の関係を直接に扱った先行研究として挙げられる。近年でも、根津茂『日本仏教を変えた親鸞の独自性：『教行信証』と『選択集』の比較から見えてきた、念仏の真価』が、法然の思想と親鸞独自の展開の境界を問うている\cite{nezu2024}。したがって、本稿は両書の比較そのものを初めて提示するものではない。

これらの先行研究は、おおまかに三つの論点に分けられる。第一に、『選択集』と『教行信証』の直接比較を通じて、専修念仏・選択本願の継承と展開を問う研究である。第二に、行信、本願力回向、真仮・方便など、親鸞側の体系化を教学的に読む研究である。第三に、『教行信証』の引用文、古写経本、引用元の扱いを検討する引用・文献関係研究である。本稿は、これら三つの論点を置き換えるものではなく、どの本文領域を精査すべきかを、本文全体の分布と近傍構造から提示する。

本稿の目的は、これらの研究を置き換えることではなく、従来研究が論じてきた法然から親鸞への継承、展開、体系化の問題が、SAT漢文本文のチャンク分布、最近傍非自己、祖師文献アンカー、語群、典拠マーカーにおいてどのように現れるかを探索的に可視化することである。先行研究が主として思想内容、教義構造、引用・文献関係、法難史や後序の文脈から両書の関係を論じてきたのに対し、本稿は、法然・親鸞・浄土教祖師文献を同一のチャンク埋め込み空間に置き、共有核、法然側に相対的に残る選択論証、親鸞側に広がる信巻・化身土巻・祖師文献方向の展開を、本文分布として観察する。

したがって、本稿は法然・親鸞比較の問題を新たに発見するものではなく、すでに蓄積された比較研究の問いを、本文全体の分布、近傍構造、三層指標によって再配置する試みである。

なお、本稿でいう「親鸞側の展開」は、法然の教えを未完成なものとして親鸞が補完したという意味ではない。本稿では、価値判断としての発展ではなく、同じ念仏圏が『選択集』では念仏を諸行から選び出す論証として現れ、『教行信証』では本願・信・証・真仏土／化身土、祖師文献引用、真仮・方便の整理を含む文類構造として展開する、という本文配置上の差に注目する。

\begin{table}[H]
\centering
\small
\caption{先行研究の関心と本稿で可視化する観点}
\label{tab:prior-research-positioning}
\begin{tabular}{@{}p{0.27\linewidth}p{0.30\linewidth}p{0.34\linewidth}@{}}
\toprule
先行研究の主な関心 & 典型的な問い & 本稿で可視化する観点 \\
\midrule
『選択集』と『教行信証』の直接比較 & 専修念仏・選択本願を親鸞がどう受けとめたか & 法然・親鸞チャンクの共有核と最近傍非自己 \\
継承と展開 & 法然の基本問題が親鸞でどう展開するか & 法然側はみ出し領域と親鸞側はみ出し領域の比較 \\
行信・本願力回向 & 念仏・行・信の位置づけはどう変わるか & 語群カウント、信巻・行巻の近接ゾーン \\
祖師文献の引用・典拠 & 親鸞はどの祖師文献をどう用いるか & 龍樹・天親・曇鸞・道綽・善導・源信アンカーへの近接 \\
真仮・方便・化身土 & 親鸞側の体系化はどこに現れるか & 化身土巻のはみ出し、道綽・源信近接、外教・方便語群 \\
\bottomrule
\end{tabular}
\end{table}

\subsection{研究上の問い}

本稿の問いは次の三点である。

\begin{enumerate}
  \item 法然『選擇本願念佛集』と親鸞『顯淨土眞實教行證文類』は、SAT漢文ベースの意味埋め込み空間でどの程度重なり、どの方向に分かれるか。
  \item 親鸞の広がりは、道綽、善導、源信、曇鸞、天親、龍樹などの浄土教祖師文献とどのような近傍関係を持つか。
  \item 共有核からのはみ出し領域を比較したとき、法然は念仏選択の論証へ、親鸞は信巻・化身土巻を中心とする典拠的体系へ分かれるのか。
\end{enumerate}

本稿の結論は、教義史的主張ではなく、意味埋め込みによって精査対象となる領域を抽出するための探索的結果である。本稿の考察は、埋め込み空間上の近傍構造、語群カウント、典拠マーカーから得られる探索的所見に限定する。既存の歴史研究、引用研究、思想史研究との全面的な照合は、本手法の信頼性や限界を測るうえで重要であるが、本稿では代表的な比較研究を確認したうえで、本分析がそれらを置き換えるものではなく、精査対象となる領域を抽出する探索地図であることを明確にする。

\section{データ}

\subsection{対象文献}

中心対象は、法然『選擇本願念佛集』と親鸞『顯淨土眞實教行證文類』である。『選擇本願念佛集』の基本的位置づけについては、浄土宗系の辞典項目および解説も参照した\cite{jodoshuzensho-senchaku,jodoshu-senchaku}。分析本文はいずれもSAT大正新脩大藏經テキストデータベースから取得した漢文本文を用いた。比較アンカーとして、浄土教祖師文献から龍樹、天親、曇鸞、道綽、善導、源信のテキストを加えた。本稿では法然を中心対象として扱うため、親鸞における七高僧のうち法然を除く祖師文献を比較アンカーとして置く。

\begin{table}[H]
\centering
\small
\caption{対象文献とチャンク数}
\label{tab:corpus}
\begin{tabular}{@{}lllllr@{}}
\toprule
分類 & 著者 & 文献 & SAT ID・範囲 & 扱い & チャンク数 \\
\midrule
中心 & 法然 & 選擇本願念佛集 & T2608 & 全体 & 69 \\
中心 & 親鸞 & 顯淨土眞實教行證文類 & T2646 & 全体 & 191 \\
アンカー & 龍樹 & 十住毘婆沙論 易行品 & T1521, 0038a25--0040a22 & 抜粋 & 6 \\
アンカー & 天親 & 無量寿経優波提舎願生偈 & T1524 & 全体 & 9 \\
アンカー & 曇鸞 & 往生論註 & T1819 & 全体 & 31 \\
アンカー & 道綽 & 安楽集 & T1958 & 全体 & 42 \\
アンカー & 善導 & 観無量寿仏経疏 & T1753 & 全体 & 23 \\
アンカー & 源信 & 往生要集 & T2682 & 全体 & 73 \\
\bottomrule
\end{tabular}
\end{table}

本稿の祖師文献アンカーは、親鸞教学における七高僧全体を網羅的に代表するものではなく、SAT上で処理可能な代表的テキストまたは抜粋を用いた探索的アンカーである。とくに龍樹は『十住毘婆沙論』易行品の抜粋6チャンク、天親は9チャンクであり、源信73チャンクなどと比べて候補数が少ない。また、善導についても本稿では『観無量寿仏経疏』を用いており、善導の他文献を含めた場合には配置が変わる可能性がある。

\subsection{本文整備と品質確認}

本文取得では、SAT HTMLから本文行を抽出し、UI部品、HTMLタグ、画像プレースホルダ、SAT行番号の本文混入を除去した。抽出本文の品質チェックでは、法然T2608が1644 SAT行、空白除去後27246文字、親鸞T2646が4215 SAT行、空白除去後76030文字であり、\texttt{U+FFFD}、HTMLタグ、HTML entity、画像プレースホルダ、ボタン文字列、SAT行番号混入はいずれも0であった。

前稿で用いた固定長トークン列を単純に文字列へ復元する方法について、本稿では比較検討を行った。その結果、チャンク境界で \texttt{U+FFFD} が発生する場合がある一方、旧方式とUnicode-safe方式の同一位置チャンクの埋め込み類似度は高く、著者別重心もほぼ一致し、大局的な意味地図への影響は小さいことを確認した。したがって前稿の広域的な配置をただちに否定するものではない。ただし、本稿では本文品質、再現性、SAT行範囲との対応をより確実にするため、トークンのオフセット情報を用い、Unicode文字境界に合わせてチャンク本文を復元する方法を採用した。本稿では、トークン境界によってUnicode文字列が破壊されないよう、Unicode文字境界を保ってチャンク本文を復元する処理をUnicode-safe方式と呼ぶ。700トークン・100トークン重なりという大枠は前稿に合わせつつ、以後の基準処理としてUnicode-safe方式を用いる。

\section{方法}

\subsection{チャンク化と埋め込み}

本文は \texttt{tiktoken} の \texttt{cl100k\_base} でトークン化し、700トークン、100トークン重なりでチャンク化した\cite{tiktoken}。チャンク本文はUnicode文字境界に合わせて作成し、各チャンクにSAT行範囲、チャンク番号、SHA-256 hashを付した。

埋め込みにはOpenAI \texttt{text-embedding-3-large} を用いた\cite{openai-embeddings}。本稿で分析した各チャンクベクトルは3072次元である。中心文献である法然・親鸞チャンクでPCAを適合し、この3072次元空間を2次元へ投影した。祖師文献アンカーは、同じ3072次元空間から固定PCA面へ投影した。これにより、法然・親鸞の配置を基準にして、各祖師文献がどの領域へ重なるかを可視化した。PCAの第1軸・第2軸の寄与率はそれぞれ0.068562、0.047405であり、2軸合計は約11.6\%である。したがって、2次元図は3072次元のうち一部を表示する探索用の投影であり、最近傍集計やコサイン類似度は元の3072次元空間で計算したものと併読する。

なお、3072次元から2次元への投影方法はPCAに限られない。UMAP、t-SNE、MDSなどを用いれば、局所近傍やクラスタの見え方が変わる可能性がある。本稿では、軸が線形結合として定義され、同じPCA面へ祖師文献アンカーを投影しやすく、再実行時の挙動も比較的安定しているため、PCAを基準表示として採用した。別の投影法による頑健性確認は、方法上の課題として残る。

また、二つの文献群が最も重なるように2次元へ投影することも、目的を明示すれば意味を持つ。たとえば法然・親鸞の共通核を見たい場合、両者を分ける方向を弱め、共有される近傍構造を強調する投影を補助図として作ることが考えられる。ただし、そのような投影は差異が見えにくくなるよう設計された図であり、重なったこと自体を独立の証拠として扱うことはできない。したがって本稿では、中立的な基準表示としてPCAを置き、共通核強調図や差異強調図は別の目的付き可視化として扱う。

\subsection{三層分析の基本構成}

本稿では、三層分析を次の基本構成として実装した。

\begin{enumerate}
  \item 意味層: \texttt{text-embedding-3-large} の3072次元コサイン類似度、2次元PCA上の距離、著者別重心、最近傍非自己、最近傍アンカー。
  \item 文体語彙層: 辞書ベースの語群カウント。語群には、念仏・称名、選択・本願、正雑・諸行、廃立・取捨、信・三心、願・回向、名号・真実、往生・浄土、方便・化身土、罪・救済を置いた。
  \item 典拠マーカー層: 経名、論釈名、祖師名、引用導入句の辞書カウント。
\end{enumerate}

ここでいう文体語彙層は、前稿で提示した広義の「文体・語彙層」を受けた名称である。ただし、本稿の実装は、厳密な文体計量や訳者判定ではなく、辞書ベースの教義語群カウントによって語彙・主題の局所的濃淡を可視化する代替指標である。

これは本格的な引用検出ではない。とくに『教行信証』は大量の引用文から構成されるため、ある語や経名が出たことは、親鸞自身の主張と引用元の主張をただちに区別しない。したがって、本稿の典拠マーカー層は、引用・参照関係そのものではなく、精査対象となる領域を抽出する探索的な代替指標である。

\begin{table}[H]
\centering
\footnotesize
\caption{文体語彙層で用いた語群辞書の概要}
\label{tab:style-lexicon}
\begin{tabular}{@{}p{0.19\linewidth}p{0.70\linewidth}@{}}
\toprule
語群 & 代表語 \\
\midrule
念仏・称名 & 念佛、念仏、稱名、称名、專念、専念、念佛三昧 \\
選択・本願 & 選擇、選択、撰擇、本願、弘誓、誓願、第十八願 \\
正雑・諸行 & 正行、雜行、雑行、二行、諸行、萬行、万行、餘行、正定業 \\
廃立・取捨 & 廢、廃、立、捨、閣、抛、傍、取、選取 \\
信・三心 & 信、信樂、信楽、三心、至誠心、深心、眞實心、一心 \\
願・回向 & 願、回向、廻向、往相、還相、利他、欲生 \\
名号・真実 & 名號、名号、眞實、真実、金剛、大利、無上 \\
往生・浄土 & 往生、淨土、浄土、安樂、安楽、極樂、極楽、報土、佛土 \\
方便・化身土 & 方便、化身土、眞門、真門、假、仮、邪、外道 \\
罪・救済 & 五逆、謗法、闡提、阿闍世、提婆、慚愧、地獄、煩惱 \\
\bottomrule
\end{tabular}
\end{table}

\begin{table}[H]
\centering
\footnotesize
\caption{典拠マーカー層で用いた明示マーカー辞書の概要}
\label{tab:source-marker-lexicon}
\begin{tabular}{@{}p{0.20\linewidth}p{0.26\linewidth}p{0.42\linewidth}@{}}
\toprule
マーカー群 & 代表語 & カウント方法と限界 \\
\midrule
経名 & 大經、觀經、阿彌陀經、無量壽經、涅槃經、大集經、日藏經、月藏經など & チャンク内の出現数を数える。異体字・新旧字体の一部は辞書に併記する。 \\
論釈名 & 往生論、往生論註、安樂集、觀經疏、法事讃、般舟讃、往生要集 & 明示的な書名・論釈名だけを数える。略称や暗示的引用は取りこぼす。 \\
祖師名 & 龍樹、天親、曇鸞、道綽、善導、源信、源空、法然、光明寺 & 人名の明示出現を数える。引用文中か地文中かは区別しない。 \\
引用導入 & 云、曰、言、釋、疏、論曰、經言、問曰、答曰、私云 & 出現数を数える。引用導入以外の一般用法を含む可能性がある。 \\
\bottomrule
\end{tabular}
\end{table}

語群・マーカーはいずれも出現数として数え、チャンク内の複数出現は重複して加算した。これは局所的濃淡を見るための単純な処理であり、語義、否定、引用者、引用元、地文との区別は行わない。

\subsection{最近傍集計と上限20サンプリング}

各チャンクについて、対象著者以外の全チャンクとの3072次元コサイン類似度を計算し、最大値を与える著者群を最近傍非自己とした。ただし、この単純最近傍集計は参照群のチャンク数に影響されうる。たとえば親鸞は191チャンク、源信は73チャンク、道綽は42チャンクであり、候補数が多い群ほど最近傍を取りやすくなる可能性がある。

そこで補助確認として、各参照群から反復ごとに最大20チャンクを抽出し、20件未満の群は全件を用いたうえで、最近傍集計を1000回再計算した。上限20サンプリング後の比率は、各反復における最近傍グループ比率の平均と95\%範囲として示す。この処理は完全な補正ではなく、候補数の多い群が単純最近傍で有利になる程度を確認する感度分析である。上限20という値は、最小群を過度に削らず、大きな群の候補数を抑えるための便宜的設定であり、上限10や30などを含む本格的な感度分析は後続の課題とする。

\subsection{はみ出しスコア}

法然側・親鸞側のはみ出し候補は、2次元PCA上の孤立度と高次元空間での非自己近傍の弱さを組み合わせて抽出した。チャンク \(i\) について、2次元PCA上で最も近い非自己チャンクまでの距離を \(d_i\)、3072次元空間で最も近い非自己チャンクとのコサイン類似度を \(s_i\) とし、次をはみ出しスコア \(p_i\) とした。

\[
  p_i = d_i (1 - s_i)
\]

このスコアは、表示面上で他群から離れ、かつ高次元空間でも非自己近傍が相対的に弱いチャンクを上位に出すための探索指標である。ただし、\(d_i\) はPCA 2次元面に依存し、本稿ではPCA座標上の距離を標準化せずに用いた。したがって、このスコアの絶対値には独立した意味を持たせず、上位候補を抽出するための順位指標として用いる。PCA 2軸の寄与率は約11.6\%にすぎないため、はみ出し表は高次元構造や歴史的関係の証明ではなく、精査対象となる領域の提示として扱う。補助的に \(1-s_i\) のみで高次元孤立度ランキングも作成し、法然側では付属・証誠・選択総結、親鸞側では信巻・化身土巻が中心になるという大枠が変わらないことを確認した。

\subsection{巻別・章別ラベル}

親鸞チャンクには、SAT行範囲に基づき、序、教巻、行巻、信巻、証巻、真仏土巻、化身土巻のラベルを付けた。法然チャンクには、序、聖道門・浄土門、正雑二行、本願念仏、三輩・一向専念、三心・信心、付属・証誠・選択総結、末尾・奥書系のラベルを付けた。

ただし、チャンクは100トークンの重なりを持つため、巻境界や章境界をまたぐことがある。ラベルはチャンク開始行に基づく近似であり、厳密な章節単位ではない。

\section{結果}

\subsection{法然・親鸞の共有核}

図\ref{fig:sat-safe-focus-map}に、SAT漢文ベースの法然・親鸞のみの意味地図を示す。PCAは法然・親鸞チャンクのみで適合し、表示範囲も両者の分布に合わせて調整した。

\begin{figure}[H]
  \centering
  \includegraphics[width=\linewidth]{../figures/sat-safe-honen-shinran-focus-map.png}
  \caption{SAT漢文・Unicode-safeチャンクによる法然・親鸞の意味地図。法然・親鸞チャンクでPCAを適合し、法然69チャンク、親鸞191チャンクを表示した。表示範囲と楕円定義は図\ref{fig:sat-safe-anchor-map}の法然・親鸞楕円と同じである。}
  \label{fig:sat-safe-focus-map}
\end{figure}

この図では、法然と親鸞は大きく重なるが、親鸞はより広い分布を示す。法然の分布は比較的まとまり、親鸞は共有核から右方向・下方向を含めて広がる。まずこの二者の関係を基準に置くことで、次に重ねる祖師文献アンカーがどの領域に入るかを読みやすくなる。

\subsection{祖師文献アンカーと親鸞の広がり}

図\ref{fig:sat-safe-anchor-map}に、同じPCA面へ龍樹、天親、曇鸞、道綽、善導、源信を投影した図を示す。PCAの適合対象は図\ref{fig:sat-safe-focus-map}と同じく法然・親鸞チャンクのみであり、祖師文献は同一の射影面へ配置している。法然・親鸞の楕円は、図\ref{fig:sat-safe-focus-map}と同じ表示範囲、同じ共分散楕円の半径定義で描いている。

また、図\ref{fig:sat-safe-focus-map}と図\ref{fig:sat-safe-anchor-map}の法然・親鸞楕円は、同一の座標、表示範囲、共分散楕円定義、描画条件で示した。祖師文献の追加は、祖師側の点群と楕円を同じ面に重ねるだけであり、法然・親鸞楕円の形状や傾きは再推定しない。

\begin{figure}[H]
  \centering
  \includegraphics[width=\linewidth]{../figures/sat-safe-honen-shinran-high-priest-anchor-map.png}
  \caption{SAT漢文・Unicode-safeチャンクによる法然・親鸞・祖師文献アンカー地図。法然・親鸞チャンクでPCAを適合し、龍樹、天親、曇鸞、道綽、善導、源信を同じ面へ投影した。法然・親鸞の表示範囲と楕円定義は図\ref{fig:sat-safe-focus-map}と同一である。}
  \label{fig:sat-safe-anchor-map}
\end{figure}

祖師文献を重ねると、道綽・善導は共有核に近く、源信・曇鸞は親鸞側の広がりと重なる。2次元表示に加え、元の3072次元空間で最近傍を集計することで、法然・親鸞間の近傍関係を確認する。

\subsection{2次元方向の後付け解釈}

PCA軸は自動的に教義ラベルを持たない。また、軸の符号は再計算によって反転しうる。したがって、図上の上下左右は、それ自体として固定した意味を持つのではなく、端にあるチャンク群の巻・章、語群、典拠マーカー、上位語から後付けで読む必要がある。ここでは、法然・親鸞チャンクの各座標について15\%/85\%分位の外側を集計し、表\ref{tab:pca-direction-reading}のように概略ラベルを付けた。

\begin{table}[H]
\centering
\small
\caption{PCA端点集計による2次元方向の概略ラベル}
\label{tab:pca-direction-reading}
\begin{tabular}{@{}llll@{}}
\toprule
方向 & 主な著者・巻/章 & 多い語群 & 概略ラベル \\
\midrule
右 / PC1+ & 親鸞、信巻・化身土巻 & 罪・救済、廃立・取捨 & 罪救済・化身土側の親鸞展開 \\
左 / PC1- & 法然、正雑二行・一向専念・三心 & 信・三心、往生・浄土、正雑・諸行 & 選択論証・共有核 \\
上 / PC2+ & 親鸞、行巻・真仏土巻 & 願・回向、往生・浄土、念仏・称名 & 阿弥陀・光明・名号・願回向 \\
下 / PC2- & 親鸞、信巻・化身土巻 & 信・三心、罪・救済、方便・化身土 & 信/三心・罪救済・方便/外教 \\
\bottomrule
\end{tabular}
\end{table}

この読みでは、横軸はおおまかに、左の「法然を多く含む選択論証・共有核」から、右の「信巻・化身土巻を中心とする親鸞側の展開」へ向かう。縦軸は、上の「阿弥陀・光明・名号・願回向」方向と、下の「信/三心・罪救済・方便/外教」方向を分けるように見える。ただし、これはこのPCA面での端点集計にもとづく記述であり、高次元空間の全構造や歴史的関係をそのまま表すものではない。

\begin{table}[H]
\centering
\small
\caption{共有核とはみ出し領域の要約}
\label{tab:shared-protrusion-summary}
\begin{tabular}{@{}p{0.17\linewidth}p{0.29\linewidth}p{0.26\linewidth}p{0.20\linewidth}@{}}
\toprule
領域 & 計算上のサイン & 主な語群・巻/章 & 解釈 \\
\midrule
共有核 & 法然・親鸞チャンクが相互に近接し、道綽・善導などのアンカーも重なる & 念仏・称名、往生・浄土、信・三心 & 専修念仏・浄土教的共通問題圏 \\
法然はみ出し & 法然チャンクが非法然点から相対的に離れ、章ラベルが選択論証部に集中する & 正雑二行、三輩・一向専念、本願念仏、付属・証誠 & 念仏を諸行から選び出す論証 \\
親鸞はみ出し & 親鸞チャンクが非親鸞点から相対的に離れ、信巻・化身土巻に集中する & 信巻、化身土巻、罪・救済、方便・化身土、廃立・取捨 & 本願・信・証・真仏土/化身土の典拠的体系 \\
\bottomrule
\end{tabular}
\end{table}

表\ref{tab:shared-protrusion-summary}は、本稿の結果を読むための要約である。以後の表と図は、この三分割を、意味層、文体語彙層、典拠マーカー層から確認するためのものである。

\begin{table}[H]
\centering
\small
\caption{法然チャンクの最近傍非自己}
\label{tab:honen-nearest}
\begin{tabular}{@{}lr@{}}
\toprule
最近傍グループ & チャンク数 \\
\midrule
親鸞 & 50 \\
道綽 & 11 \\
源信 & 4 \\
善導 & 3 \\
曇鸞 & 1 \\
\bottomrule
\end{tabular}
\end{table}

\begin{table}[H]
\centering
\small
\caption{親鸞チャンクの最近傍非自己}
\label{tab:shinran-nearest}
\begin{tabular}{@{}lr@{}}
\toprule
最近傍グループ & チャンク数 \\
\midrule
法然 & 48 \\
道綽 & 45 \\
源信 & 36 \\
曇鸞 & 23 \\
天親 & 21 \\
善導 & 17 \\
龍樹 & 1 \\
\bottomrule
\end{tabular}
\end{table}

\begin{figure}[H]
  \centering
  \includegraphics[width=\linewidth]{../figures/shared-core-protrusion-nearest-bars.png}
  \caption{共有核とはみ出し領域と最近傍非自己の可視化。上段は表\ref{tab:shared-protrusion-summary}の定性的要約を横流れ図にしたもので、幅は量を表さない。下段は表\ref{tab:honen-nearest}・表\ref{tab:shinran-nearest}の単純最近傍非自己数を横棒で示す。}
  \label{fig:shared-protrusion-bars}
\end{figure}

表\ref{tab:honen-nearest}では、法然チャンク69件のうち50件が親鸞を最近傍非自己とする。この結果は、少なくとも本分析の本文・チャンク・埋め込み条件では、法然チャンクの多くが、親鸞文献チャンクと意味空間上で強く重なることを示す。ただし、埋め込みの近さそのものは影響関係やその向きを証明しない。本稿では影響史の説明へ踏み込まず、意味空間上の近傍関係として扱う。

表\ref{tab:shinran-nearest}では、親鸞チャンク191件の最近傍非自己が法然だけでなく、道綽、源信、曇鸞、天親、善導へ広く分かれる。図\ref{fig:shared-protrusion-bars}の下段で見ると、法然側は親鸞への集中が目立つのに対し、親鸞側は法然・道綽・源信を中心に広く分散する。親鸞は法然との共有核を持つ一方で、祖師文献群の複数方向へ展開している。

\begin{table}[H]
\centering
\small
\caption{参照群上限20サンプリングによる最近傍比率}
\label{tab:nearest-subsampling}
\begin{tabular}{@{}llrp{0.15\linewidth}p{0.29\linewidth}@{}}
\toprule
対象 & 最近傍群 & 全件/抽出数 & 単純比率 & 上限20サンプリング平均比率（95\%範囲） \\
\midrule
法然 & 親鸞 & 191/20 & 72.5\% & 33.9\%（11.6--55.1\%） \\
法然 & 道綽 & 42/20 & 15.9\% & 27.2\%（11.6--43.5\%） \\
法然 & 善導 & 23/20 & 4.3\% & 18.2\%（10.1--27.5\%） \\
法然 & 源信 & 73/20 & 5.8\% & 15.7\%（2.9--31.9\%） \\
親鸞 & 道綽 & 42/20 & 23.6\% & 22.8\%（15.2--29.3\%） \\
親鸞 & 法然 & 69/20 & 25.1\% & 19.0\%（11.5--26.2\%） \\
親鸞 & 天親 & 9/9 & 11.0\% & 15.9\%（13.1--19.4\%） \\
親鸞 & 曇鸞 & 31/20 & 12.0\% & 14.4\%（9.4--19.4\%） \\
\bottomrule
\end{tabular}
\end{table}

表\ref{tab:nearest-subsampling}は、参照群の候補数を上限20件に制限した1000回反復の感度分析である。法然から親鸞への単純最近傍比率は72.5\%と高いが、親鸞群を20件に制限すると平均33.9\%へ下がる。したがって、単純最近傍には親鸞チャンク数の多さによる影響が含まれる。一方で、上限20サンプリング後も法然側では親鸞・道綽が上位に残り、親鸞側では道綽・法然・天親・曇鸞へ分散する。このため、本稿では単純最近傍を探索指標とし、上限20サンプリング結果を、共有核とはみ出しを読む際の注意として併読する。

\subsection{親鸞の巻別・近接ゾーン}

親鸞チャンクを巻別に分け、最近傍非自己と2次元上の孤立度から、法然近接、祖師近接、親鸞側突出候補に分類した。ここでいう親鸞側突出候補は、教学史的な意味での親鸞独自思想を直接判定するものではなく、本分析条件において非親鸞チャンクとの近傍が相対的に弱い箇所を指す。表\ref{tab:shinran-volume-affinity}に主要な巻別・近接ゾーンを示す。

\begin{table}[H]
\centering
\small
\caption{親鸞の巻別・近接ゾーン上位}
\label{tab:shinran-volume-affinity}
\begin{tabular}{@{}llr@{}}
\toprule
巻 & 近接ゾーン & チャンク数 \\
\midrule
信巻 & 法然近接 & 21 \\
化身土巻 & 道綽近接 & 19 \\
行巻 & 法然近接 & 14 \\
信巻 & 親鸞側突出候補 & 13 \\
証巻 & 天親近接 & 11 \\
化身土巻 & 法然近接 & 10 \\
化身土巻 & 源信近接 & 9 \\
信巻 & 源信近接 & 8 \\
真仏土巻 & 善導近接 & 8 \\
行巻 & 源信近接 & 8 \\
\bottomrule
\end{tabular}
\end{table}

信巻には法然近接と親鸞側突出候補がともに多く現れる。これは、法然から受けた専修念仏の枠組みが、親鸞において信、三心、真実心、回向、本願力の問題へ展開することと対応しうる。化身土巻では道綽近接、源信近接、善導近接、親鸞側突出候補が混在する。これは、方便、末法、外教批判、護法・鬼神、真仮の整理などが、単一の祖師文献方向ではなく複合的に広がることを示す。

\subsection{法然三層チャンク列: 選択論証への収束}

以下では、共有核とはみ出し領域を、本文順の三層チャンク列として確認する。これにより、法然側の選択論証がどの位置で濃く現れるか、また親鸞側の信巻・化身土巻がどの祖師文献方向・語群・典拠マーカーと重なるかを、本文順に見ることができる。

図\ref{fig:honen-three-layer}に、法然『選擇本願念佛集』69チャンクを横軸に取り、意味層、文体語彙層、典拠マーカー層を縦に並べた三層ヒートマップを示す。意味層は、各法然チャンクから親鸞および祖師文献各群への3072次元コサイン類似度の最大値である。文体語彙層と典拠マーカー層は、親鸞図と同じ辞書ベースの語群・明示マーカーを用いた。

\begin{figure}[H]
  \centering
  \includegraphics[width=\linewidth]{../figures/honen-three-layer-sequence-heatmap.png}
  \caption{法然『選択集』三層チャンク列。上段は意味層、中段は文体語彙層、下段は典拠マーカー層である。横軸は法然チャンク順、縦線はSAT行範囲ベースの章・位置境界を示す。文体語彙層は語群カウントの代替指標であり、典拠マーカー層は引用検出そのものではない。}
  \label{fig:honen-three-layer}
\end{figure}

この図では、意味層の多くが親鸞文献チャンクと近く出る一方、章・位置によって道綽、善導、源信へ寄る箇所も分かれる。文体語彙層では、正雑二行、本願念仏、三輩・一向専念、三心・信心、付属・証誠・選択総結に、念仏・称名、選択・本願、正雑・諸行、信・三心、往生・浄土の語群が帯状に現れる。典拠マーカー層では引用導入句が全体に強く、七高僧名や経名の出方だけでは章構成を十分に分けられない。したがって法然側の三層図からは、法然チャンクが親鸞文献チャンクと意味的に重なりつつ、『選択集』内部では念仏を諸行から選び出す論証へ収束していく様子が読み取れる。

\subsection{親鸞三層チャンク列: 典拠体系への拡張}

図\ref{fig:shinran-three-layer}に、親鸞『教行信証』191チャンクを横軸に取り、意味層、文体語彙層、典拠マーカー層を縦に並べた三層ヒートマップを示す。意味層は、各親鸞チャンクから法然および祖師文献各群への3072次元コサイン類似度の最大値である。文体語彙層は、辞書ベースの語群カウントである。典拠マーカー層は、浄土三部経、涅槃経、大集/日月蔵、論/論註、安楽集、観経疏/善導釈、往生要集、七高僧名、引用導入句の明示マーカーである。

\begin{figure}[H]
  \centering
  \includegraphics[width=\linewidth]{../figures/shinran-three-layer-sequence-heatmap.png}
  \caption{親鸞『教行信証』三層チャンク列。上段は意味層、中段は文体語彙層、下段は典拠マーカー層である。横軸は親鸞チャンク順、縦線は巻境界を示す。文体語彙層は語群カウントの代替指標であり、典拠マーカー層は引用検出そのものではない。}
  \label{fig:shinran-three-layer}
\end{figure}

この図では、意味層と文体語彙層が必ずしも同じ動きをしない。たとえば信巻では、意味層の最近傍は法然、源信、道綽に分かれる一方、文体語彙層では信・三心と罪・救済が強く出る。化身土巻では、意味層で道綽、法然、源信が混在し、文体語彙層では信・三心、願・回向、往生・浄土、方便・化身土、廃立・取捨が分散する。典拠マーカー層では引用導入句が全巻を通じて強く、特定典拠名だけでは『教行信証』の引用構造を十分に分けられないことも分かる。したがって、この三層図からは、意味的近さ、語彙的特徴、明示的引用マーカーが一致する箇所と対応しない箇所を区別できる。

\subsection{法然のはみ出し領域}

法然チャンクのうち、親鸞・祖師文献アンカーから相対的に離れる領域を抽出した。抽出には、2次元PCA上の最近傍非法然距離と、高次元コサイン類似度を組み合わせたはみ出しスコア（protrusion score）を用いた。表\ref{tab:honen-protrusion}には上位10件を示し、本文では補助集計として上位24件の章分布を用いる。

\begin{table}[H]
\centering
\small
\caption{法然はみ出し領域の上位チャンク}
\label{tab:honen-protrusion}
\begin{tabular}{@{}rlllll@{}}
\toprule
番号 & 位置 & 行開始 & 近傍 & 近傍祖師 & 語群 \\
\midrule
7 & 正雑二行 & T2608, 83.0003a04 & 親鸞 & 善導 & 正雑・諸行 \\
44 & 付属・証誠・選択総結 & T2608, 83.0013b17 & 親鸞 & 曇鸞 & 念仏・称名 \\
22 & 三輩・一向専念 & T2608, 83.0007b07 & 親鸞 & 善導 & 念仏・称名 \\
29 & 三心・信心 & T2608, 83.0009b05 & 親鸞 & 曇鸞 & 念仏・称名 \\
59 & 付属・証誠・選択総結 & T2608, 83.0017c18 & 道綽 & 道綽 & 念仏・称名 \\
56 & 付属・証誠・選択総結 & T2608, 83.0017a01 & 親鸞 & 道綽 & 念仏・称名 \\
15 & 本願念仏 & T2608, 83.0005b05 & 親鸞 & 源信 & 往生・浄土 \\
37 & 三心・信心 & T2608, 83.0011b21 & 親鸞 & 源信 & 信・三心 \\
21 & 三輩・一向専念 & T2608, 83.0007a11 & 道綽 & 道綽 & 正雑・諸行 \\
55 & 付属・証誠・選択総結 & T2608, 83.0016b22 & 親鸞 & 道綽 & 往生・浄土 \\
\bottomrule
\end{tabular}
\end{table}

法然のはみ出し領域には、正雑二行、三輩・一向専念、本願念仏、三心・信心、付属・証誠・選択総結が多い。これは、念仏そのものの単純な称揚というより、諸行の中から念仏を選び出し、専修の行として位置づける論証の密度を反映していると考えられる。

\subsection{親鸞のはみ出し領域}

同じ尺度で、親鸞チャンクについても非親鸞点からの2次元距離と高次元コサイン類似度を組み合わせ、親鸞側のはみ出し領域を抽出した。表\ref{tab:shinran-protrusion}には上位10件を示し、本文では補助集計として上位24件の巻・内容ゾーン分布を用いる。

\begin{table}[H]
\centering
\small
\caption{親鸞はみ出し領域の上位チャンク}
\label{tab:shinran-protrusion}
\begin{tabular}{@{}rlllll@{}}
\toprule
番号 & 巻 & 内容ゾーン & 行開始 & 近傍 & 語群 \\
\midrule
162 & 化身土巻 & 護法・鬼神・宇宙秩序 & T2646, 83.0635b21 & 源信 & 往生・浄土 \\
84 & 信巻 & 信/三心・罪救済 & T2646, 83.0614a08 & 源信 & 罪・救済 \\
72 & 信巻 & 信/三心・罪救済 & T2646, 83.0610b22 & 源信 & 罪・救済 \\
83 & 信巻 & 信/三心・罪救済 & T2646, 83.0613c14 & 道綽 & 罪・救済 \\
85 & 信巻 & 信/三心・罪救済 & T2646, 83.0614b02 & 源信 & 廃立・取捨 \\
71 & 信巻 & 信/三心・罪救済 & T2646, 83.0610a27 & 源信 & 願・回向 \\
182 & 化身土巻 & 外教批判 & T2646, 83.0640b28 & 道綽 & 廃立・取捨 \\
75 & 信巻 & 信/三心・罪救済 & T2646, 83.0611b10 & 道綽 & 罪・救済 \\
74 & 信巻 & 信/三心・罪救済 & T2646, 83.0611a10 & 源信 & 罪・救済 \\
78 & 信巻 & 信/三心・罪救済 & T2646, 83.0612b05 & 源信 & 罪・救済 \\
\bottomrule
\end{tabular}
\end{table}

上位24件を集計すると、巻別には信巻13件、化身土巻8件、真仏土巻2件、証巻1件であった。内容ゾーンでは、信巻の信・三心・罪救済が13件、化身土巻の護法・鬼神・宇宙秩序が4件、真仏土巻の名号・本願・真実が2件、化身土巻の外教批判が1件、方便・真仮整理が1件であった。最近傍非自己は源信10件、道綽6件、曇鸞4件、法然3件、善導1件であり、親鸞のはみ出しは法然から単に離れるというより、源信・道綽・曇鸞などの引用・主題圏と接しながら、信巻と化身土巻の問題群へ広がる。

\subsection{法然・親鸞のはみ出し比較}

法然はみ出し上位24件は、付属・証誠・選択総結10件、三輩・一向専念4件、本願念仏4件、三心・信心3件、正雑二行3件であった。これに対して親鸞は、信巻13件、化身土巻8件を中心に広がる。したがって、同じ「共有核からのはみ出し」でも、法然側は念仏を諸行から選び出す論証の密度として現れ、親鸞側は信巻における罪と救済、化身土巻における方便・真仮・外教批判・護法秩序のような、体系化の周縁部として現れる。

三層チャンク列で見ても、この違いはおおむね対応する。法然側では、意味層で親鸞との近さを保ちつつ、文体語彙層が選択・本願、正雑・諸行、念仏・称名、往生・浄土へまとまる。親鸞側では、意味層が法然、源信、道綽などへ分かれ、文体語彙層も信・三心、罪・救済、方便・化身土、廃立・取捨へ広がる。つまり、法然の三層図は論証の収束を、親鸞の三層図は典拠的体系の拡張を示している。

\subsection{代表箇所の確認}

以上のはみ出し領域が、実際の本文読解へ戻る際にどのような精査候補になるかを、本文を転載せず、SAT行範囲、近傍、語群、マーカーから3例だけ確認する。ここでの確認は、原文の歴史的読解を代替するものではなく、計算上抽出された箇所がどの読解方向を示すかを点検するためのものである。以下の3例は、法然側の選択論証、親鸞信巻の罪・救済、親鸞化身土巻の外教・方便整理という、本稿で抽出された三つの主要な突出領域を代表する箇所として選んだ。いずれも、はみ出しスコア上位10件または上位24件集計で繰り返し現れた主要ゾーンに属する。

\begin{table}[H]
\centering
\footnotesize
\setlength{\tabcolsep}{2pt}
\caption{代表箇所の確認}
\label{tab:case-study-chunks}
\begin{tabular}{@{}p{0.09\linewidth}p{0.17\linewidth}p{0.20\linewidth}p{0.16\linewidth}p{0.25\linewidth}@{}}
\toprule
対象 & 箇所 & SAT行範囲 & 近傍・語群 & 読解上の確認点 \\
\midrule
法然 & chunk 7, 正雑二行 & T2608, 83.0003a04--0003b04 & 最近傍は親鸞、近傍祖師は善導、語群は正雑・諸行 & 正行・雑行の分類を通じて、念仏を諸行から選び出す論証として読める。 \\
親鸞 & chunk 84, 信巻 & T2646, 83.0614a08--0614b05 & 最近傍は源信、語群は罪・救済 & 信巻の罪・救済問題が、法然から単に離れるのではなく、源信方向の主題圏と接しながら現れる。 \\
親鸞 & chunk 182, 化身土巻 & T2646, 83.0640b28--0640c27 & 最近傍は道綽、語群は廃立・取捨 & 化身土巻の外教批判・真仮整理が、道綽方向の問題圏と接しつつ、親鸞側の体系化の周縁部として現れる。 \\
\bottomrule
\end{tabular}
\end{table}

第一の例では、法然側のはみ出しが、念仏という主題だけでなく、正行・雑行を分ける分類操作として現れている。最近傍非自己は親鸞であり、意味的には共有核に近いが、語群は正雑・諸行に集中する。この組み合わせが、「同じ念仏圏にいながら、法然側では選択論証が濃く出る」という本稿の読みを支える。

第二の例では、親鸞側のはみ出しが信巻に集中し、罪・救済の語群を伴う。最近傍が源信であることは、親鸞の信巻が法然との共有核から離れて孤立するというより、往生要集的な救済・罪業の問題圏と接しながら、信・三心の議論へ広がることを示す精査候補になる。

第三の例では、化身土巻の外教批判・真仮整理が、道綽に近い点として抽出される。ここでは、典拠マーカーだけで引用関係を断定することはできないが、意味層と語群層がともに、真仮・方便・外教批判の周縁的問題を示す。化身土巻の護法・鬼神・宇宙秩序に関わる別の上位候補（chunk 162）も源信方向へ近く出ており、化身土巻のはみ出しは一つの典拠方向ではなく、複数の祖師文献方向に接しながら広がる。

\subsection{はみ出しスコアの頑健性確認}

はみ出しスコアは2次元PCA距離を含むため、PCA面に依存する。そこで補助的に、2次元距離を使わず、元の3072次元空間での高次元孤立度 \(1-s_i\) だけによる上位24件を全チャンクから抽出し、はみ出しスコア上位24件と比較した。表\ref{tab:high-dim-isolation-check}に結果を示す。

\begin{table}[H]
\centering
\footnotesize
\setlength{\tabcolsep}{2pt}
\caption{はみ出しスコアと高次元孤立度の上位24件比較}
\label{tab:high-dim-isolation-check}
\begin{tabular}{@{}p{0.08\linewidth}p{0.25\linewidth}p{0.25\linewidth}p{0.12\linewidth}p{0.20\linewidth}@{}}
\toprule
対象 & はみ出し上位24件の中心 & 高次元孤立度上位24件の中心 & 重なり & 解釈 \\
\midrule
法然 & 付属・証誠・選択総結10、三輩・一向専念4、本願念仏4、正雑二行3、三心・信心3 & 付属・証誠・選択総結8、三心・信心5、正雑二行4、本願念仏2、聖道門・浄土門2 & 11件、J=0.297 & 個別順位は変わるが、選択論証・付属証誠・正雑二行は残る。 \\
親鸞 & 信巻13、化身土巻8、真仏土巻2、証巻1 & 化身土巻13、信巻7、真仏土巻3、証巻1 & 11件、J=0.297 & 順位は入れ替わるが、信巻と化身土巻中心という大枠は残る。 \\
\bottomrule
\end{tabular}
\end{table}

この比較では、上位集合の完全一致は見られない。これは、はみ出しスコアが2次元PCA上の孤立を含むのに対し、高次元孤立度は元空間の最近傍類似度だけを見るためである。ただし、法然側では付属・証誠・選択総結、正雑二行、本願念仏が残り、親鸞側では信巻・化身土巻への集中が残る。したがって、本稿の大きな読みは2次元PCA距離だけに依存していないが、個別チャンク順位は探索指標として扱う必要がある。

\section{考察}

\subsection{解釈の範囲}

以下の考察は、埋め込み空間上の配置、高次元コサイン類似度、2次元PCA上のはみ出し、辞書ベースの語群、典拠マーカーから得た探索的所見である。本稿の結果は、既存研究を置き換える結論ではなく、既存研究と照合すべき精査候補を抽出する探索地図である。将来的には、こうした外部研究との照合によって、本手法がどの程度信頼できるか、どこで誤読や過剰解釈を生むかを検証する必要がある。本稿はその前段階として、方法そのものがどのような領域を抽出するかを記述する。

\subsection{法然は「念仏を選ぶ論理」を構成する}

法然『選択集』の特徴は、念仏を諸行から選び出す論証にある。本分析で抽出したはみ出し領域でも、正雑二行、諸行との取捨、本願、三輩・一向専念、付属・証誠が強く現れた。したがって、法然の独自領域は、本分析の範囲では「念仏の意味が親鸞と異なる」というより、「なぜ諸行ではなく念仏なのか」を構成する分類・廃立の論証として読むのが自然である。

この読みは、『選択集』という書物の形式とも合う。『選択集』は、浄土三部経や善導などを引きながら、称名念仏を往生行として選び出す書である。親鸞と共有する内容は多いが、書物としての焦点は比較的鋭く、念仏選択の論理へ収束する。

\subsection{親鸞は「本願・信・証の典拠的体系」へ展開する}

親鸞『教行信証』は、本分析の意味地図では、法然に近い領域を持ちながら、道綽、源信、曇鸞、天親、善導へ広がった。とくに信巻や化身土巻では、法然近接と親鸞側突出候補、祖師近接が混在する。

このことから、親鸞は単に法然を反復するのではなく、法然の念仏選択を、本願、信、回向、名号、真仏土、化身土、方便、罪と救済、外教批判などを含む大きな典拠的体系へ再配置していると読める。とくに親鸞側のはみ出し領域が信巻と化身土巻に集中したことは、この再配置が中心教義の反復だけではなく、救済可能性の限界事例、方便と真実の区別、外教批判や護法秩序まで含む体系拡張として現れていることを示唆する。

\subsection{三層分析の意義}

ただし、意味層だけでは解釈が混ざる。たとえば、親鸞のあるチャンクが曇鸞に近いとしても、それは曇鸞の引用文が含まれるからなのか、親鸞自身の論述が曇鸞的主題を展開しているからなのか、現行の処理では区別できない。同様に、法然のはみ出し領域に善導や道綽が近い場合も、実際の引用、語彙、主題、書物構成が混在している。

したがって、意味層、文体語彙層、典拠マーカー層を分けることには解釈上の意味がある。意味層は主題の近さを示す。文体語彙層は、選択、正雑、廃立、信、願、回向、名号、真実などの語彙操作を示す。典拠マーカー層は、経名、人名、引用導入句、固有句によって、どの文献群を明示的に参照しているかを示す。この三層が一致する箇所と対応しない箇所を区別することで、意味的近接をただちに影響関係や教義的一致として読んでしまう危険を抑えることができる。

既存の『選択集』・『教行信証』比較研究は、専修念仏、選択本願、行信、本願力回向、真仮・方便などの思想史的・教学的問題を論じてきた。本稿の分析は、それらの解釈を直接検証するものではないが、法然側の選択論証領域と、親鸞側の信巻・化身土巻・祖師文献方向の広がりを、本文分布上の精査候補として提示する。したがって、本稿の三層探索地図は、既存研究を置き換えるものではなく、既存研究に戻って読むべき箇所を選び出す補助線として位置づけられる。

\section{限界}

本稿には少なくとも次の限界がある。

\begin{enumerate}
  \item 2次元PCAは全分散の約11.6\%のみを示すため、図上の距離は高次元空間の近傍関係と必ずしも一致しない。
  \item 2次元投影法はPCAに限られない。UMAP、t-SNE、MDSなどで同じ傾向が出るかは未検証である。
  \item チャンクは700トークン・100トークン重なりであり、巻・章・引用単位などの自然境界とは一致しない。
  \item 巻別・章別ラベルはSAT行範囲の開始位置に基づく近似であり、境界をまたぐチャンクがある。
  \item 語群カウントは辞書ベースの限定的な実装であり、否定、引用者、文脈、語義の違いを見ていない。
  \item 典拠マーカー層は引用検出そのものではない。引用文の立場と著者自身の立場を分ける処理が必要である。『教行信証』の引用文を古写経本から再検討する研究\cite{fujiwara2020}との照合も、今後の重要な検証課題である。
  \item 本稿の考察は埋め込み空間上の近傍構造、語群カウント、典拠マーカーから得られる探索的所見であり、個別箇所について既存研究と全面的に照合した結論ではない。
  \item 古写経本、異本、写本系資料はまだ直接分析に組み込んでいない。それらの研究成果との照合は、本手法の信頼性や限界を検証する後続の課題である。
\end{enumerate}

以上の制約を踏まえると、引用文と地文の分離、法然『選択集』の十六章門との対応、親鸞『教行信証』の巻別・引用別ラベルの精密化が、次の検証課題となる。また、法然側の比較資料を広げる場合には、『黒谷上人語灯録』系\cite{jodoshuzensho-gotoroku}や三部経釈系を、本文品質と伝来上の注意を確認したうえで扱う必要がある。

\section{結論}

本稿では、法然『選択集』、親鸞『教行信証』、および浄土教祖師文献アンカーを、SAT漢文・Unicode-safeチャンク・\texttt{text-embedding-3-large} によって三層探索地図化した。法然と親鸞は意味空間上で強く重なる。しかし、本稿で重要なのは、重なるかどうかではなく、共有核から何がどの方向へはみ出すかである。

法然側のはみ出し領域は、念仏理解の単純な差というより、念仏を諸行から選び出す論証、すなわち正雑二行、取捨、選択、本願、付属・証誠のまとまりとして現れる。親鸞側のはみ出し領域は、信巻の罪と救済、化身土巻の方便・真仮整理、護法・鬼神・宇宙秩序、外教批判として現れ、法然側とは異なる広がり方を示した。最近傍の上限20サンプリングを行うと、単純最近傍に含まれるチャンク数バイアスは低下するが、法然側の親鸞・道綽近接と、親鸞側の道綽・法然・天親・曇鸞への分散という大きな読みは残る。

したがって、本分析の範囲では、法然と親鸞の差は、念仏そのものの意味差としてよりも、同じ念仏圏をどのような論証形式・典拠体系へ展開するかの差として現れる。法然側では、正雑二行、取捨、選択、本願、付属・証誠を通じて念仏を諸行から選び出す論証が相対的に濃く現れる。一方、親鸞側では、信巻の罪と救済、化身土巻の方便・真仮整理、護法・鬼神・宇宙秩序、外教批判が、道綽・源信・曇鸞・天親などの祖師文献方向への広がりとともに現れる。これは既存研究を置き換える結論ではなく、歴史研究・引用研究との比較によって検証すべき精査候補である。

従来の『選択集』・『教行信証』比較研究は、専修念仏・選択本願の継承、行信・本願力回向、真仮・方便、祖師文献引用などの論点から、法然と親鸞の関係を論じてきた。本稿の結果は、これらの論点を新たに結論づけるものではないが、法然側の選択論証への収束と、親鸞側の信巻・化身土巻・祖師文献方向への分散を、本文全体の近傍構造として可視化する。したがって、本稿は既存研究の結論を置き換えるものではなく、既存研究が精査してきた問題領域を、本文分布上の候補領域として再提示するものである。

\begin{thebibliography}{99}

\bibitem{ueyama2026three-layer-map}
上山大信「意味・文体・典拠マーカーの三層地図による仏教文献探索：宗派別参照群・阿弥陀経二訳・親鸞文献の予備的分析」2026.
\url{https://dueyama.github.io/buddhist-text-embedding-map/paper/}
（2026年6月5日閲覧）

\bibitem{sat}
SAT大正新脩大藏經テキストデータベース.
\url{https://21dzk.l.u-tokyo.ac.jp/SAT/}
（2026年6月5日閲覧）

\bibitem{moriwaki1953}
森脇一掬「選択集と教行信証に関する一考察」『龍谷論叢』1, 48--72, 1953.
INBUDS ID: IB00020630A.
\url{https://tripitaka.l.u-tokyo.ac.jp/INBUDS/}
（2026年6月5日閲覧）

\bibitem{inaba1972}
稲葉秀賢「『教行信証』と『選択集』」『親鸞教学』21, 33--50, 1972.
\url{https://otani.repo.nii.ac.jp/records/10720}
（2026年6月5日閲覧）

\bibitem{umehara1974}
梅原隆章「選択集と教行信証」『日本文化と浄土教論攷：井川定慶博士喜寿記念』362--371, 1974.
INBUDS ID: IB00046839A.
\url{https://tripitaka.l.u-tokyo.ac.jp/INBUDS/}
（2026年6月5日閲覧）

\bibitem{asai2004}
浅井成海「『選択集』と『教行信証』――基本的問題の継承と展開――」『行信学報』18, 2004.
\url{https://www.gyousin.com/institute/journal/}
（2026年6月5日閲覧）

\bibitem{nezu2024}
根津茂『日本仏教を変えた親鸞の独自性：『教行信証』と『選択集』の比較から見えてきた、念仏の真価』法藏館, 2024.
\url{https://ndlsearch.ndl.go.jp/books/R100000001-I2711B17026510}
（2026年6月5日閲覧）

\bibitem{jodoshuzensho-senchaku}
新纂浄土宗大辞典「選択本願念仏集」.
\url{https://jodoshuzensho.jp/daijiten/index.php/%E9%81%B8%E6%8A%9E%E6%9C%AC%E9%A1%98%E5%BF%B5%E4%BB%8F%E9%9B%86}
（2026年6月5日閲覧）

\bibitem{jodoshuzensho-gotoroku}
新纂浄土宗大辞典「黒谷上人語灯録」.
\url{https://jodoshuzensho.jp/daijiten/index.php/%E9%BB%92%E8%B0%B7%E4%B8%8A%E4%BA%BA%E8%AA%9E%E7%81%AF%E9%8C%B2}
（2026年6月5日閲覧）

\bibitem{jodoshu-senchaku}
浄土宗公式サイト「法然上人のご生涯と浄土宗の教え 選択集」.
\url{https://jodo.or.jp/newspaper/special/6226/}
（2026年6月5日閲覧）

\bibitem{fujiwara2020}
藤原智「『教行信証』における引用文について：古写経本による再検討」『近現代『教行信証』研究検証プロジェクト研究紀要』3, 63--107, 2020.
\url{https://www.jstage.jst.go.jp/article/kyogyoshinsho/3/0/3_63/_article/-char/ja/}
（2026年6月5日閲覧）

\bibitem{openai-embeddings}
OpenAI, Embeddings documentation.
\url{https://platform.openai.com/docs/guides/embeddings}
（2026年6月5日閲覧）

\bibitem{tiktoken}
OpenAI, \texttt{tiktoken}.
\url{https://github.com/openai/tiktoken}
（2026年6月5日閲覧）

\end{thebibliography}

\end{document}
